\documentclass[10pt,a4paper]{amsart}
\usepackage[margin=19mm]{geometry}
\usepackage{amsmath,amssymb,mathtools}
\usepackage[scr=rsfs,cal=euler]{mathalfa}
\usepackage{xfrac}
\usepackage[unicode,hidelinks,bookmarks=false]{hyperref}
\newcommand{\ie}{i.e.\ }
\newcommand{\cf}{cf.\ }
\newcommand{\resp}{respectively\ }
\newcommand{\Subsection}{\subsection}
\providecommand{\nobreakdash}{\nobreak\hbox{-}}
\setlength{\emergencystretch}{2em}
\multlinegap0pt
\usepackage{amssymb}
\usepackage{tikz}
\newtheorem{theorem}{Theorem}
\title{Defect of irreducible plane curves with simple singularities}
\author{Piotr Pokora}
\date{Source: arXiv:2404.03341v2 (CC BY 4.0)}
\begin{document}
\maketitle

\noindent\emph{Source and licence.} Defect of irreducible plane curves with simple singularities. Version arXiv:2404.03341v2 (CC BY 4.0), distributed by arXiv under the Creative Commons Attribution 4.0 International licence, http://creativecommons.org/licenses/by/4.0/, as recorded in the arXiv licence metadata for this exact version and shown on the abstract page https://arxiv.org/abs/2404.03341v2. Full licence notice: \texttt{licenses/arxiv-2404.03341-CC-BY-4.0.txt}.\par\smallskip

\noindent\textit{Editorial scope.} Counted unit: the article's theorem deff (source0.tex lines 119--127). The article's complete original proof of it is delivered with it (source0.tex lines 199--213, one \texttt{proof} environment of 15 lines, which the article places away from the statement). No mathematics is written, repaired or re-derived: the statement and the proof are the article's own lines, in the article's own notation, and no retained line is substituted or re-flowed.  Retained uncounted as the source-local context the proof needs: lines 180--184.  Nothing was omitted from any retained span, so the package carries no omission record.  The retained lines cite 2 works, whose entries are transcribed verbatim from the article's own bibliography source in the e-print and delivered with short editorial labels recorded in the item's reference mapping.  No retained line contains a figure environment, an image-inclusion command or drawing code, so no image is delivered and the package declares no asset dependency.  Editorial formatting only, in addition: an AMS \texttt{amsart} preamble that restates the article's own declarations and macros used by the retained lines, namely the environments \texttt{theorem} on separate counters, together with the standard packages those lines need.  The retained environments print in this document as Theorem 1 in order of appearance, whereas the article prints the same items with the numbering of its own full document.  The complete native source of the article as distributed in the arXiv e-print for this exact version is bundled unchanged as \texttt{sources/157206/source0.tex}.
\begin{theorem}[Dimca-Sernesi]
\label{sern}
Let $C \, = \, \{ f=0 \}$ be a reduced curve of degree $d$ in $\mathbb{P}^{2}_{\mathbb{C}}$ having only quasi-homogeneous singularities. Then $${\rm mdr}(f) \geq \alpha_{C}\cdot d - 2,$$
where $\alpha_{C}$ denotes the Arnold exponent of $C$.
\end{theorem}

\begin{theorem}[{\cite[Theorem 1.2]{Dimca1}}]
\label{deff}
Let $C \, = \, \{ f=0 \}$ be a reduced plane curve of degree $d$ and $r= {\rm mdr}(C)$. Then the following hold.
\begin{itemize}
    \item If $r < (d-1)/2$, then $\nu(C) = (d-1)^{2}-r(d-1-r)-\tau(C)$.
    \item If $r \geq (d-2)/2$, then
    $$\nu(C) = \bigg\lceil\frac{3}{4}(d-1)^{2}\bigg\rceil - \tau(C).$$
\end{itemize}
\end{theorem}

\begin{proof}
 The existence of nodal curves is granted by a result due to Severi \cite{Severi}. Since all singular points $p \in {\rm Sing}(C_{d})$ are nodes, we have ${\rm lct}_{p}(C_{d}) =1$, and the Arnold exponent of $C_{d}$ is equal to
 $$\alpha_{C_{d}} = 1.$$
 Then by Theorem \ref{sern} we have
 $${\rm mdr}(C_{d}) \geq d-2.$$
 By the assumption $d\geq 4$, so the following inequality holds
 $$d - 2 \geq \frac{d-2}{2},$$
 which means that by Theorem \ref{deff} the defect of $C_{d}$ is equal to 
 $$\nu(C_{d}) = \bigg\lceil\frac{3}{4}(d-1)^{2}\bigg\rceil - \tau(C_{d}).$$
 Now we want to find an upper bound on $\tau(C_{d})$. First of all, since all singularities of $C_{d}$ are nodes, one has $\tau_{p}(C_{d}) = 1$ for every $p \in {\rm Sing}(C_{d})$. Since the number of nodes of $C_{d}$ is bounded from above by $\frac{(d-1)(d-2)}{2}$, we get
 $$\tau(C_{d}) \leq \frac{(d-1)(d-2)}{2}.$$
 Taking into account the above inequality, we finally get 
 $$\nu(C_{d}) \geq \frac{3}{4}(d-1)^{2} - \frac{(d-1)(d-2)}{2} = \frac{1}{4}(d^{2}-1),$$
 which completes the proof.
\end{proof}

\begin{thebibliography}{99}
\bibitem[Dimca1]{Dimca1} A. Dimca, On rational cuspidal plane curves and the local cohomology of Jacobian rings. \textit{Comment. Math. Helv.} \textbf{94(4)}: 689 -- 700 (2019).
\bibitem[Severi]{Severi} F. Severi, \textit{Vorlesungen \"uber algebraische Geometrie}. Teubner (1921).
\end{thebibliography}
\end{document}
